% References to labels of other documents whose \externaldocument gives
% xr-hyper's URL of that document's PDF, loaded after hyperref: after the
% file, on the next line, after a space, empty, and holding ~ and _; with
% the nocite option, through \externalcitedocument, and with no prefix.
% The page builds, and each reference prints as without a URL: from an
% .aux that LaTeX wrote, its number, page, or name as in the PDF, with no
% link, in each context; from an .aux that tex4ht wrote, with its .xref, a
% link to that document's page.
\documentclass{article}
\usepackage{amsthm}
\usepackage{hyperref}
\usepackage{xr-hyper}
\externaldocument[X-]{test-xr/latex}[https://example.org/ext.pdf]
\externaldocument[H-][nocite]{test-xrurl/hyper}
  [https://example.org/hyper.pdf]
\externalcitedocument{test-xrurl/tex4ht}[https://example.org/~a/b_c.pdf]
\externaldocument[P-]{test-xr/plain} []
\newtheorem{theorem}{Theorem}
\begin{document}
\section{See \ref{X-ss}}
CaseA \ref{X-ss}, page \pageref{X-front}, named \nameref{X-thm},
auto \autoref{X-thm}.

CaseB \ref{H-oss}, auto \autoref{H-oss}.

CaseC \ref{oss}.

CaseD \ref{P-ss}, page \pageref{P-front}.

\textit{In italics, \ref{X-ss}.}
A note.\footnote{In a footnote, \ref{X-ss}.}
\textsf{In sans-serif, \ref{X-ss}.}
\texttt{In typewriter, \ref{X-ss}.}

\begin{theorem}
In a theorem, \ref{X-ss}.
\end{theorem}

\section{Here}\label{here}
Local \ref{here}.
\end{document}
